% Deutsche Version
\documentclass[12pt,a4paper]{article}
\usepackage[utf8]{inputenc}
\usepackage[T1]{fontenc}
\usepackage{lmodern}
\usepackage{amsmath,amssymb}
\usepackage{graphicx}
\usepackage{xcolor}
\usepackage{hyperref}
\usepackage{geometry}
\geometry{a4paper, left=3cm, right=3cm, top=3cm, bottom=3cm}
\usepackage{setspace}
\onehalfspacing
\usepackage{parskip}
\usepackage[ngerman]{babel}
\usepackage{csquotes}
\usepackage{microtype}
\usepackage{booktabs}
\usepackage{longtable}
\usepackage{array}
\usepackage{listings}
\usepackage{caption}
\usepackage{subcaption}
\usepackage{float}
\usepackage{url}
\usepackage{natbib}
\usepackage{titling}

\lstset{
  language=Python,
  basicstyle=\ttfamily\small,
  keywordstyle=\color{blue},
  commentstyle=\color{green!40!black},
  stringstyle=\color{red},
  showstringspaces=false,
  numbers=left,
  numberstyle=\tiny,
  numbersep=5pt,
  breaklines=true,
  frame=single,
  backgroundcolor=\color{gray!5},
  tabsize=2,
  captionpos=b
}

\title{\Huge\textbf{ARS als Metamethodologie} \\[2mm]
       \LARGE Die Bedingungen erklärender Modelle \\[2mm]
       \LARGE im Zeitalter generativer KI}

\author{
  \large
  \begin{tabular}{c}
    Paul Koop
  \end{tabular}
}

\date{\large 1994--2026}

\begin{document}

\maketitle

\begin{abstract}
Dieser Beitrag betrachtet die Algorithmisch Rekursive Sequenzanalyse (ARS) nicht 
nur als Methode, sondern als \textit{Metamethodologie} – ein Rahmenwerk, das die 
Bedingungen spezifiziert, unter denen ein Modell als \textit{erklärend} und nicht 
nur als \textit{beschreibend} oder \textit{simulierend} gelten kann. Ausgehend 
von den Kernprinzipien der ARS – dem Primat der Interpretation, der Trennung von 
Struktur und Statistik, der kontrollierten Falsifikation und der XAI-Validierung 
– argumentiere ich, dass diese Prinzipien notwendige Kriterien für erklärende 
Modelle in jeder Disziplin darstellen, die sich mit sequenziellen sozialen 
Prozessen befasst. Der Beitrag setzt die ARS systematisch zu fünf zeitgenössischen 
Forschungsprogrammen in Beziehung: (1) Formale Verifikation und Model Checking, 
(2) Interpretierbares maschinelles Lernen und Rule Extraction, (3) Grounded 
Theory, (4) Kausale Inferenz und (5) Process Mining. Für jedes dieser Felder 
zeige ich, was die ARS von diesem Ansatz lernen kann und, entscheidend, was 
dieser Ansatz von der ARS lernen kann. Der Beitrag schließt mit der These, dass 
die ARS einen transdisziplinären Benchmark dafür bereitstellt, echte Erklärung 
von statistischer oder struktureller Beschreibung zu unterscheiden.
\end{abstract}

\newpage
\tableofcontents
\newpage

\section{Einleitung: Von der Methode zur Metamethodologie}

Die Algorithmisch Rekursive Sequenzanalyse (ARS) wurde in ihren verschiedenen 
Versionen (2.0–4.0) primär als eine Methode zur formalen Analyse sequenzieller 
Interaktionen präsentiert. Sie transformiert interpretativ gewonnene Kategorien 
in formale Modelle: probabilistische kontextfreie Grammatiken (PCFG), Petri-Netze, 
Bayessche Netze und deterministische endliche Automaten (DFA). Dieser Beitrag 
nimmt eine andere Perspektive ein. Er argumentiert, dass die ARS nicht nur eine 
Methode ist, sondern eine \textit{Metamethodologie} – ein Rahmenwerk, das die 
\textit{Bedingungen} spezifiziert, unter denen ein Modell als erklärend gelten 
kann.

Diese Perspektivverschiebung wird durch eine wiederkehrende Verwirrung in der 
zeitgenössischen KI und Datenwissenschaft motiviert: die Gleichsetzung von 
\textbf{Simulation} mit \textbf{Erklärung}. Ein großes Sprachmodell (LLM), das 
auf Verkaufsgesprächen trainiert wurde, kann plausible Dialoge mit hoher 
Genauigkeit simulieren. Doch wie die ARS-Notizbücher gezeigt haben, kann es 
nicht erklären, \textit{warum} eine bestimmte Sequenz von Sprechakten wohlgeformt 
ist oder welche Regeln ihre Generierung konstituieren. Das LLM liefert einen 
statistischen Schatten des Prozesses; die ARS liefert sein strukturelles Gerüst.

Die metamethodologische These dieses Beitrags ist, dass die ARS-Prinzipien 
\textbf{notwendige Bedingungen} für erklärende Modelle in jeder Disziplin 
darstellen, die sich mit sequenziellen sozialen Prozessen befasst. Diese 
Bedingungen sind:

\begin{enumerate}
    \item \textbf{Interpretative Verankerung}: Symbole müssen an dokumentierte 
    qualitative Interpretationen gebunden sein, nicht nur an statistische 
    Korrelationen.
    
    \item \textbf{Strukturelle Entscheidbarkeit}: Die Wohlgeformtheit von 
    Sequenzen muss formal entscheidbar sein (z.B. durch einen DFA), unabhängig 
    von empirischen Häufigkeiten.
    
    \item \textbf{Generative Transparenz}: Das Modell muss in der Lage sein, 
    Sequenzen so zu generieren, dass jeder Schritt auf explizite Regeln 
    zurückgeführt werden kann.
    
    \item \textbf{Falsifizierbarkeit}: Interpretationen müssen der kontrollierten 
    Falsifikation unterliegen; Gegenbeispiele müssen Regeln widerlegen können.
    
    \item \textbf{XAI-Validierung}: Das Modell muss die NIST-XAI-Kriterien 
    erfüllen – Verständlichkeit, Genauigkeit, Wissensgrenzen.
\end{enumerate}

Um diese These zu untermauern, setze ich die ARS systematisch zu fünf zeitgenössischen 
Forschungsprogrammen in Beziehung, die überlappende Anliegen teilen: formale 
Verifikation, interpretierbares maschinelles Lernen, Grounded Theory, kausale 
Inferenz und Process Mining. Für jedes stelle ich zwei reziproke Fragen:

\begin{enumerate}
    \item Was kann die ARS von diesem Ansatz lernen? (Technische oder konzeptionelle 
    Verbesserungen)
    
    \item Was kann dieser Ansatz von der ARS lernen? (Methodologische Sicherungen, 
    Kriterien für Erklärung)
\end{enumerate}

\section{ARS als Metamethodologie: Die Kernprinzipien}

Bevor die fünf Ansätze untersucht werden, müssen die ARS-Prinzipien dargelegt 
werden, die als metamethodologischer Benchmark dienen.

\subsection{Prinzip 1: Interpretative Verankerung}

In der ARS ist jedes Terminalzeichen (KBG, VBBd, KAA etc.) das Produkt einer 
dokumentierten qualitativen Interpretation. Der Kodierungsprozess ist keine 
Black Box; er wird aufgezeichnet, begründet und unterliegt der intersubjektiven 
Validierung. Dieses Prinzip schließt rein datengetriebene Kategorienbildung 
(z.B. Clustering von Embeddings) als Erklärung aus.

\subsection{Prinzip 2: Trennung von Struktur und Statistik}

Wie in \texttt{ARS\_XAI\_Aut2\_Ger.tex} entwickelt, hält die ARS eine strikte 
Trennung zwischen strukturellen Regeln (die deterministisch und kontextfrei 
sind) und statistischen Regularitäten (die empirisch und kontingent sind) ein. 
Eine strukturelle Regel gilt oder gilt nicht – unabhängig davon, wie oft sie 
verletzt wird. Diese Trennung fehlt in rein statistischen Modellen.

\subsection{Prinzip 3: Generative Transparenz}

Die induzierte Grammatik muss in der Lage sein, Sequenzen auf nachvollziehbare 
Weise zu generieren. Der Transduktor in Lisp, der Parser in Pascal und der 
Induktor in Scheme bieten jeweils ein unterschiedliches Fenster in diese 
Transparenz. Ein Modell, das keine Exemplare aus seinen eigenen Regeln 
generieren kann, ist nicht erklärend.

\subsection{Prinzip 4: Kontrollierte Falsifikation}

Interpretationen werden nicht ein für alle Mal behauptet. Sie werden als 
Lesarten produziert und dann durch nachfolgende Sequenzpositionen falsifiziert 
(nach Oevermanns Sequenzanalyse). Dies erzeugt eine Spirale von Interpretation 
und Widerlegung, die Poppers Falsifikationismus für qualitatives Material 
adaptiert.

\subsection{Prinzip 5: XAI-Validierung}

Die drei NIST-XAI-Kriterien – Verständlichkeit, Genauigkeit, Wissensgrenzen – 
sind keine optionalen Zusätze, sondern konstitutive Elemente erklärender 
Modelle. Verständlichkeit erfordert semantische Interpretierbarkeit; Genauigkeit 
erfordert Übereinstimmung mit dem Material; Wissensgrenzen erfordern explizite 
Dokumentation der Modellgrenzen.

\section{Fünf Ansätze im Dialog mit der ARS}

\subsection{Formale Verifikation und Model Checking}

\subsubsection{Was formale Verifikation ist}

Formale Verifikation, insbesondere Model Checking, ist eine Methode aus der 
theoretischen Informatik, die systematisch und algorithmisch prüft, ob ein 
formales Modell (z.B. ein endlicher Automat, ein Petri-Netz oder ein Bayessches 
Netz) spezifizierte Eigenschaften erfüllt. Diese Eigenschaften umfassen 
\textit{Safety} ("etwas Schlechtes passiert nie") und \textit{Liveness} 
("etwas Gutes passiert irgendwann").

Model Checking ist exhaustiv: Es erkundet den gesamten Zustandsraum des Modells. 
Im Gegensatz zum statistischen Testen, das probabilistische Garantien liefert, 
bietet Model Checking \textit{deterministische} Garantien über das Verhalten 
des Modells.

\subsubsection{Was die ARS von formaler Verifikation lernen kann}

\begin{itemize}
    \item \textbf{Spezifikationssprachen für Eigenschaften}: Die ARS könnte 
    temporale Logiken (LTL, CTL) übernehmen, um zu spezifizieren, welche 
    Eigenschaften die induzierte Grammatik erfüllen sollte. Zum Beispiel: 
    $\square (KBG \rightarrow \lozenge VBG)$ ("Immer, wenn ein Kunde grüßt, 
    grüßt irgendwann der Verkäufer zurück").
    
    \item \textbf{Generierung von Gegenbeispielen}: Wenn eine Eigenschaft 
    verletzt wird, produzieren Model Checker eine Gegenbeispielspur. Dies 
    könnte als ein mächtiges Falsifikationswerkzeug für ARS-Interpretationen 
    dienen.
    
    \item \textbf{Bewusstsein für Zustandsraumexplosion}: ARS-Modellierer 
    sollten sich bewusst sein, dass hierarchische Grammatiken zu großen 
    Zustandsräumen führen können. Formale Verifikation bietet Abstraktionstechniken 
    zur Bewältigung dieses Problems.
\end{itemize}

\subsubsection{Was formale Verifikation von der ARS lernen kann}

\begin{itemize}
    \item \textbf{Interpretative Verankerung von Zuständen}: Im standardmäßigen 
    Model Checking sind Zustände abstrakte Symbole. Die ARS besteht darauf, 
    dass jeder Zustand interpretativ verankert sein muss. Dies könnte zu einem 
    neuen Teilgebiet führen: \textit{interpretatives Model Checking}, bei dem 
    Eigenschaften geprüft \textit{und} die Bedeutung der Zustände dokumentiert 
    wird.
    
    \item \textbf{Die Trennung von Struktur und Statistik}: Model Checking 
    geht typischerweise von einem deterministischen Modell aus. Die ARS zeigt, 
    wie strukturelle Regeln (die verifiziert werden können) von statistischen 
    Variationen (die es nicht können) getrennt werden können. Dies könnte das 
    probabilistische Model Checking bereichern.
    
    \item \textbf{XAI-Kriterien für Verifikation}: Die Ergebnisse des Model 
    Checkings sind für Nicht-Experten oft opak. Die XAI-Kriterien der ARS 
    könnten die Entwicklung verständlicherer Verifikationsausgaben leiten.
\end{itemize}

\subsection{Interpretierbares maschinelles Lernen und Rule Extraction}

\subsubsection{Was IML und Rule Extraction sind}

Interpretierbares maschinelles Lernen (IML) zielt darauf ab, Black-Box-Modelle 
(neuronale Netze, Gradient Boosting, Random Forests) für den Menschen 
verständlich zu machen. \textit{Rule Extraction} (Regel-Extraktion) ist eine 
spezifische IML-Technik, die versucht, die gelernte Funktion approximativ 
durch eine Menge von Wenn-Dann-Regeln zu beschreiben (z.B. mit RIPPER, CART 
oder durch Analyse von Aktivierungsmustern).

Im Gegensatz zur ARS, die Regeln direkt aus Daten induziert, arbeitet Rule 
Extraction typischerweise \textit{post-hoc}: Das Modell ist bereits trainiert, 
und Regeln werden als Erklärung extrahiert.

\subsubsection{Was die ARS von IML und Rule Extraction lernen kann}

\begin{itemize}
    \item \textbf{Skalierbare Regelinduktion}: Die ARS induziert derzeit 
    Regeln aus kleinen Korpora (n = 8). IML bietet Techniken zur Extraktion 
    von Regeln aus großen Datensätzen, wenn auch mit geringerer interpretativer 
    Kontrolle.
    
    \item \textbf{Quantitative Regelbewertung}: IML bietet Metriken zur 
    Bewertung extrahierter Regeln (Coverage, Fidelity, Stabilität). Die ARS 
    könnte diese übernehmen, um zu beurteilen, wie gut eine Grammatik 
    generalisiert.
    
    \item \textbf{Hybride Regelmengen}: Einige IML-Methoden kombinieren 
    globale und lokale Regeln. Die ARS könnte hybride Grammatiken mit einem 
    strukturellen Kern und lokalen statistischen Variationen erforschen.
\end{itemize}

\subsubsection{Was IML und Rule Extraction von der ARS lernen können}

\begin{itemize}
    \item \textbf{Explikation vs. Approximation}: Die Rule Extraction in IML 
    ist fast immer approximativ. Die ARS besteht auf \textit{exakten} Regeln 
    für das gegebene Korpus. Dies wirft eine grundlegende Frage auf: Ist 
    Approximation für Erklärung jemals akzeptabel? Die ARS legt nahe: 
    Approximation ist nur akzeptabel, wenn der Approximationsfehler dokumentiert 
    ist und der strukturelle Kern exakt bleibt.
    
    \item \textbf{Interpretative Validierung von Regeln}: IML extrahiert Regeln, 
    die statistisch zu den Daten passen. Die ARS fügt eine Ebene der 
    \textit{interpretativen Validierung} hinzu: Regeln müssen auch für 
    menschliche Interpreten sinnvoll sein. Dies könnte die Extraktion von 
    statistisch korrekten, aber semantisch bedeutungslosen Regeln verhindern.
    
    \item \textbf{Falsifizierbarkeit als Kriterium}: IML diskutiert selten, 
    wie extrahierte Regeln falsifiziert werden könnten. Die ARS macht 
    Falsifizierbarkeit zu einem zentralen Kriterium. IML könnte dies übernehmen: 
    Eine Regelmenge ist nicht erklärend, wenn kein denkbares Gegenbeispiel 
    sie widerlegen könnte.
\end{itemize}

\subsection{Grounded Theory}

\subsubsection{Was Grounded Theory ist}

Die Grounded Theory (GT) nach Glaser und Strauss ist eine klassische 
Methodologie der qualitativen Sozialforschung zur Entwicklung von Theorien 
aus Daten. Sie umfasst Verfahren wie offenes Kodieren, axiales Kodieren und 
selektives Kodieren. Das Ziel ist die Generierung von Theorien mittlerer 
Reichweite, die im empirischen Material "verankert" (grounded) sind.

In jüngerer Zeit gibt es Versuche, Teile der GT zu formalisieren oder 
computergestützt zu unterstützen (z.B. durch Natural Language Processing 
oder Topic Modeling). Die GT bleibt jedoch in ihrer endgültigen Ausgabe 
weitgehend informell.

\subsubsection{Was die ARS von der Grounded Theory lernen kann}

\begin{itemize}
    \item \textbf{Systematische Kodierverfahren}: Die GT bietet ein reiches 
    Vokabular und eine Reihe von Kodierverfahren, die die interpretative 
    Phase der ARS bereichern könnten. Konzepte wie "axiales Kodieren" 
    (Beziehung von Kategorien zu Unterkategorien) ähneln der hierarchischen 
    Kompression der ARS, sind aber feinkörniger.
    
    \item \textbf{Theoretisches Sampling}: Das Prinzip des theoretischen 
    Samplings – Auswahl neuer Fälle basierend auf sich entwickelnden 
    theoretischen Einsichten – könnte die Fallauswahl der ARS in größeren 
    Studien leiten.
    
    \item \textbf{Ständiger Vergleich}: Die Methode des ständigen Vergleichs 
    (Vergleich jedes neuen Falls mit bereits entwickelten Kategorien) ist 
    bereits implizit im systematischen Fallvergleich der ARS (Phase 4) 
    enthalten, könnte aber expliziter gemacht werden.
\end{itemize}

\subsubsection{Was die Grounded Theory von der ARS lernen kann}

\begin{itemize}
    \item \textbf{Von der Theorie zum generativen Modell}: Die GT stoppt 
    typischerweise auf der Ebene narrativer Theorie oder Kategoriensysteme. 
    Die ARS geht weiter: Sie transformiert die Theorie in eine \textit{generative 
    Grammatik}, die neue Sequenzen produzieren kann. Die GT könnte dies 
    übernehmen, um ihre Theorien testbar und ausführbar zu machen.
    
    \item \textbf{Formale Falsifizierbarkeit}: Die Validierungsverfahren der 
    GT sind primär qualitativ (z.B. Member Checking, Peer Debriefing). Die 
    ARS fügt formale Falsifizierbarkeit hinzu: Die Grammatik kann auf eine 
    Weise falsch sein, die mechanisch demonstriert werden kann (z.B. durch 
    einen Parser, der eine Sequenz ablehnt).
    
    \item \textbf{XAI-Kriterien für Grounded Theory}: Die XAI-Kriterien der 
    ARS (Verständlichkeit, Genauigkeit, Wissensgrenzen) bieten eine Checkliste 
    für die Bewertung grounded theories. Eine GT-Theorie, die ihre 
    Wissensgrenzen nicht spezifizieren kann, ist unvollständig.
\end{itemize}

\subsection{Kausale Inferenz und kausale Graphenmodelle}

\subsubsection{Was kausale Inferenz ist}

Kausale Inferenz, insbesondere mit kausalen Graphenmodellen (z.B. DAGs, DoWhy, 
CausalNex), geht über bloße Korrelation hinaus und versucht, kausale Beziehungen 
zwischen Variablen zu identifizieren und zu quantifizieren. Sie verwendet 
Techniken wie den Do-Kalkül, instrumentelle Variablen und kontrafaktisches 
Schließen.

Eine zentrale Einsicht der kausalen Inferenz ist, dass Korrelation nicht 
Kausalität ist. Gerichtete azyklische Graphen (DAGs) werden verwendet, um 
Annahmen über kausale Strukturen zu repräsentieren.

\subsubsection{Was die ARS von kausaler Inferenz lernen kann}

\begin{itemize}
    \item \textbf{Kausale Interpretation von Grammatiken}: ARS-Grammatiken 
    beschreiben sequenzielle Abhängigkeiten. Kausale Inferenz könnte helfen 
    zu unterscheiden, ob diese Abhängigkeiten lediglich sequenziell oder 
    genuin kausal sind. Zum Beispiel: Verursacht die Frage "Sonst noch 
    etwas?" einen zusätzlichen Kauf oder ist sie nur damit korreliert?
    
    \item \textbf{Kontrafaktisches Schließen}: Kausale Inferenz exzelliert 
    bei der Beantwortung kontrafaktischer Fragen ("Was wäre passiert, wenn 
    der Verkäufer nicht gefragt hätte?"). Die ARS könnte die kontrafaktische 
    Simulation (bereits vorhanden in Phase 3 der CGTI) als Standardvalidierungswerkzeug 
    übernehmen.
    
    \item \textbf{Instrumentelle Variablen für sequenzielle Daten}: Die ARS 
    befasst sich mit sequenziellen Daten, bei denen Confounding häufig ist. 
    Techniken mit instrumentellen Variablen könnten helfen, kausale Effekte 
    selbst in beobachteten sequenziellen Daten zu identifizieren.
\end{itemize}

\subsubsection{Was kausale Inferenz von der ARS lernen kann}

\begin{itemize}
    \item \textbf{Interpretative Verankerung kausaler Graphen}: In der 
    standardmäßigen kausalen Inferenz wird der DAG oft ohne interpretative 
    Dokumentation angenommen oder aus Daten gelernt. Die ARS besteht darauf, 
    dass jeder Knoten und jede Kante interpretativ verankert sein muss. Dies 
    könnte zu einer \textit{interpretativen kausalen Inferenz} als neuem 
    Teilgebiet führen.
    
    \item \textbf{Sequenzielle Grammatiken als kausale Strukturen}: ARS-Grammatiken 
    sind eine Form der kausalen Struktur über Sequenzen. Die kausale Inferenz 
    befasst sich typischerweise mit statischen oder Zeitreihendaten, nicht mit 
    grammatikalischen Sequenzen. Die ARS könnte eine neue Klasse \textit{grammatikalischer 
    kausaler Modelle} inspirieren.
    
    \item \textbf{XAI für kausale Modelle}: Kausale Modelle werden oft als 
    DAGs mit Wahrscheinlichkeiten präsentiert, die nicht selbsterklärend sind. 
    Die XAI-Kriterien der ARS könnten die Dokumentation kausaler Modelle 
    leiten und sie für Domänenexperten zugänglicher machen.
\end{itemize}

\subsection{Process Mining}

\subsubsection{Was Process Mining ist}

Process Mining ist ein Forschungsfeld an der Schnittstelle von Data Mining, 
maschinellem Lernen und Prozessmodellierung. Es zielt darauf ab, aus 
Ereignislogs – sequenziellen Aufzeichnungen von Prozessschritten, z.B. in 
Workflow-Management-Systemen oder ERP-Systemen – Prozessmodelle (häufig in 
Form von Petri-Netzen, BPMN-Diagrammen oder direkten Folgengraphen) zu 
entdecken, zu konformieren und zu verbessern.

Process Mining arbeitet typischerweise mit großen, anonymisierten und schwach 
annotierten Logs. Es entdeckt den "durchschnittlichen Prozess" oder die 
"häufigsten Pfade". Es zielt nicht auf eine vollständige Rekonstruktion der 
konstitutiven Regeln eines Einzelfalls ab.

\subsubsection{Was die ARS von Process Mining lernen kann}

\begin{itemize}
    \item \textbf{Skalierbare Entdeckungsalgorithmen}: Process Mining bietet 
    ausgefeilte Algorithmen (z.B. Alpha-Miner, Heuristics-Miner, Inductive 
    Miner) zur Entdeckung von Petri-Netzen aus großen Logs. Die ARS könnte 
    diese für größere Korpora übernehmen oder anpassen, während sie die 
    interpretative Kontrolle bewahrt.
    
    \item \textbf{Conformance Checking}: Process Mining umfasst Techniken 
    zur Prüfung, ob ein Ereignislog mit einem gegebenen Modell konform ist. 
    Dies könnte zur Validierung von ARS-Grammatiken an neuen Daten verwendet 
    werden.
    
    \item \textbf{Leistungsanalyse}: Process Mining fügt Leistungsdimensionen 
    (Zeit, Kosten, Häufigkeit) hinzu. Die ARS könnte erweitert werden, um 
    zeitliche und ressourcenbezogene Dimensionen systematischer zu integrieren.
\end{itemize}

\subsubsection{Was Process Mining von der ARS lernen kann}

\begin{itemize}
    \item \textbf{Interpretative Entdeckung für kleine Logs}: Process Mining 
    benötigt typischerweise große Logs, um zuverlässige Modelle zu produzieren. 
    Die ARS zeigt, wie aus einem einzigen Fall (n = 1) durch interpretative 
    Tiefe Modelle entdeckt werden können. Dies könnte für Process Mining in 
    Bereichen mit knappen Daten wertvoll sein (z.B. medizinische Verfahren, 
    juristische Fälle).
    
    \item \textbf{Dokumentation von Entdeckungsentscheidungen}: Process-Mining-Algorithmen 
    treffen viele Entscheidungen (z.B. welche Pfade einbezogen werden, wie mit 
    Rauschen umgegangen wird). Diese Entscheidungen werden selten in interpretativ 
    zugänglicher Weise dokumentiert. Die reflexive Dokumentation der ARS könnte 
    als Vorbild dienen.
    
    \item \textbf{Trennung von Struktur und Statistik}: Process Mining produziert 
    oft Modelle, die strukturelle Regeln mit statistischem Rauschen vermischen. 
    Die strikte Trennung der ARS könnte die Qualität der entdeckten Modelle 
    verbessern, indem klar unterschieden wird, was strukturell notwendig ist 
    und was nur empirisch häufig vorkommt.
    
    \item \textbf{XAI für Process Mining}: Die von Process Mining produzierten 
    Modelle (z.B. spaghettiartige Petri-Netze) sind oft schwer verständlich. 
    Die XAI-Kriterien der ARS könnten die Entwicklung verständlicherer 
    Process-Mining-Ausgaben leiten.
\end{itemize}

\section{Auf dem Weg zu einem transdisziplinären Benchmark für Erklärung}

Die fünf Vergleiche oben offenbaren ein gemeinsames Muster. Jeder zeitgenössische 
Ansatz hat technische Stärken, die die ARS verbessern könnten: Skalierbarkeit 
(IML, Process Mining), formale Rigorosität (Verifikation, kausale Inferenz) 
und systematische Kodierverfahren (Grounded Theory). Umgekehrt fehlt jedem 
Ansatz ein Teil der metamethodologischen Sicherungen, die die ARS bietet: 
interpretative Verankerung, strukturelle Entscheidbarkeit, generative 
Transparenz, kontrollierte Falsifikation und XAI-Validierung.

Diese Symmetrie legt nahe, dass die ARS nicht nur eine Methode unter anderen 
ist, sondern ein \textbf{transdisziplinärer Benchmark} dafür, was als 
erklärendes Modell gelten kann.

\begin{table}[H]
\centering
\caption{ARS als transdisziplinärer Benchmark}
\label{tab:benchmark}
\begin{tabular}{@{} p{3cm} p{4cm} p{6cm} @{}}
\toprule
\textbf{Kriterium} & \textbf{Frage} & \textbf{Abwesenheit bedeutet...} \\
\midrule
Interpretative Verankerung & Sind die Symbole sinnhaft dokumentiert? & Statistische Korrelation ohne Verstehen \\
Strukturelle Entscheidbarkeit & Ist Wohlgeformtheit formal entscheidbar? & Probabilistisches Raten statt Regelbefolgung \\
Generative Transparenz & Kann das Modell Exemplare nachvollziehbar generieren? & Simulation ohne Erklärung \\
Kontrollierte Falsifikation & Können Gegenbeispiele Regeln widerlegen? & Unfalsifizierbares post-hoc Storytelling \\
XAI-Validierung & Sind Verständlichkeit, Genauigkeit und Wissensgrenzen dokumentiert? & Technische Raffinesse ohne epistemische Rechenschaftspflicht \\
\bottomrule
\end{tabular}
\end{table}

\subsection{Die Unterscheidung von Erklärung und Simulation neu betrachtet}

Die Kernunterscheidung, die aus dieser Analyse hervorgeht, ist die zwischen 
\textbf{Erklärung} und \textbf{Simulation}. Ein Modell \textit{simuliert}, 
wenn es die statistischen Eigenschaften der Daten reproduziert. Ein Modell 
\textit{erklärt}, wenn es die konstitutiven Regeln spezifiziert, die das 
Phänomen erzeugen.

\begin{itemize}
    \item \textbf{Simulation} ist ausreichend für Prognose. Ein LLM, das das 
    nächste Token genau vorhersagt, ist ein guter Simulator.
    \item \textbf{Erklärung} ist notwendig für Verstehen, Intervention und 
    normative Bewertung. Eine ARS-Grammatik, die die Regeln eines Verkaufsgesprächs 
    spezifiziert, ist eine Erklärung.
\end{itemize}

Die fünf Kriterien oben sind die Bedingungen, unter denen ein Modell als 
Erklärung und nicht nur als Simulation qualifiziert.

\subsection{Die Rolle des menschlichen Interpreten}

Ein wiederkehrendes Thema in allen fünf Vergleichen ist die Rolle des 
menschlichen Interpreten. In der ARS ist der Mensch konstitutiv: Interpretation 
ist ein menschlicher Akt, der nicht vollständig automatisiert werden kann. In 
den anderen Ansätzen ist der Mensch oft extern – entwirft Algorithmen, stellt 
Trainingsdaten bereit, evaluiert Ausgaben.

Dies ist keine Schwäche der ARS, sondern eine Stärke. Die ARS macht explizit, 
was oft implizit bleibt: dass Erklärung eine menschliche Praxis ist, keine 
Eigenschaft eines isoliert betrachteten Modells. Ein Modell ist erklärend 
\textit{für jemanden}, der es verstehen, nutzen und dafür Rechenschaft 
ablegen kann.

\section{Fazit und Ausblick}

Dieser Beitrag hat argumentiert, dass die Algorithmisch Rekursive Sequenzanalyse 
(ARS) nicht nur eine Methode ist, sondern eine Metamethodologie – ein Rahmenwerk, 
das die Bedingungen spezifiziert, unter denen ein Modell als erklärend gelten 
kann. Fünf Kernprinzipien wurden identifiziert: interpretative Verankerung, 
strukturelle Entscheidbarkeit, generative Transparenz, kontrollierte Falsifikation 
und XAI-Validierung.

Der Beitrag setzte die ARS dann systematisch zu fünf zeitgenössischen 
Forschungsprogrammen in Beziehung: formaler Verifikation, interpretierbarem 
maschinellen Lernen, Grounded Theory, kausaler Inferenz und Process Mining. 
Für jedes wurde gezeigt, was die ARS von diesen Ansätzen lernen kann 
(technische Verbesserungen) und, entscheidend, was diese Ansätze von der 
ARS lernen können (methodologische Sicherungen, Kriterien für Erklärung).

Die metamethodologische These ist, dass die ARS-Prinzipien notwendige 
Bedingungen für erklärende Modelle in jeder Disziplin darstellen, die sich 
mit sequenziellen sozialen Prozessen befasst. Sie bieten einen transdisziplinären 
Benchmark für die Unterscheidung echter Erklärung von statistischer oder 
struktureller Beschreibung.

Für die weitere Forschung sind drei Richtungen besonders vielversprechend:

\begin{enumerate}
    \item \textbf{Implementierung hybrider Systeme}: Integration von 
    ARS-Grammatiken mit Model Checkern, Rule Extractors oder Process-Mining-Algorithmen 
    unter Wahrung der metamethodologischen Sicherungen.
    
    \item \textbf{Empirische Prüfung des Benchmarks}: Anwendung der fünf 
    Kriterien auf bestehende Modelle in verschiedenen Disziplinen und 
    Prüfung, ob sie die wahrgenommene Erklärungsqualität vorhersagen.
    
    \item \textbf{Erweiterung auf nicht-sequenzielle Domänen}: Obwohl die 
    ARS für sequenzielle Daten entwickelt wurde, könnten die metamethodologischen 
    Prinzipien auf andere Modelltypen generalisieren (z.B. Klassifikation, 
    Clustering, Regression).
\end{enumerate}

Abschließend: Die Frage ist nicht, ob ein Modell zu den Daten passt. 
Statistische Passung ist notwendig, aber nicht hinreichend. Die Frage ist, 
ob das Modell die metamethodologischen Kriterien erfüllt, die Erklärung 
erst möglich machen. Die ARS bietet eine konkrete, operationalisierte 
Antwort.

\newpage
\begin{thebibliography}{99}

\bibitem[Baier \& Katoen(2008)]{baier2008principles}
Baier, C., \& Katoen, J.-P. (2008). \textit{Principles of Model Checking}. 
MIT Press.

\bibitem[Glaser \& Strauss(1967)]{glaser1967discovery}
Glaser, B. G., \& Strauss, A. L. (1967). \textit{The Discovery of Grounded 
Theory: Strategies for Qualitative Research}. Aldine.

\bibitem[Koop(1992)]{koop1992parser}
Koop, P. (1992). \textit{Demo-Parser Chart-Parser Version 1.0}. Pascal-Quellcode.

\bibitem[Koop(1994)]{koop1994scheme}
Koop, P. (1994). \textit{Grammatikinduktion empirisch gesicherter 
Verkaufsgespräche}. Scheme-Quellcode.

\bibitem[Koop(1994)]{koop1994lisp}
Koop, P. (1994). \textit{Sequenzanalyse empirisch gesicherter 
Verkaufsgespräche}. Lisp-Quellcode.

\bibitem[Koop(2023)]{koop2023notebook}
Koop, P. (2023). \textit{Qualitative Sozialforschung und Große Sprachmodelle}. 
Jupyter Notebook.

\bibitem[Koop(2024)]{koop2024ars}
Koop, P. (2024/2026). \textit{Zwischen Interpretation und Berechnung: 
Algorithmisch Rekursive Sequenzanalyse als Brücke zwischen qualitativer 
Hermeneutik und formaler Modellierung}. the-last-freedom.org.

\bibitem[Koop(2026)]{koop2026causal}
Koop, P. (2026). \textit{Soziale Strukturen und Prozesse: Kausale Inferenz 
mit Probabilistischen kontextfreien Grammatiken und Bayesschen Netzen}. 
the-last-freedom.org.

\bibitem[Molnar(2022)]{molnar2022interpretable}
Molnar, C. (2022). \textit{Interpretable Machine Learning: A Guide for 
Making Black Box Models Explainable}. leanpub.com.

\bibitem[Oevermann et al.(1979)]{oevermann1979methodology}
Oevermann, U., Allert, T., Konau, E., \& Krambeck, J. (1979). Die Methodologie 
einer objektiven Hermeneutik. In H.-G. Soeffner (Hrsg.), \textit{Interpretative 
Verfahren in den Sozial- und Textwissenschaften} (S. 352-434). Metzler.

\bibitem[Ortigossa et al.(2024)]{ortigossa2024xai}
Ortigossa, E. S., Gonçalves, T., \& Nonato, L. G. (2024). Explainable 
Artificial Intelligence (XAI)—From Theory to Methods and Applications. 
\textit{IEEE Access}, 12, 80799-80846.

\bibitem[Pearl(2009)]{pearl2009causality}
Pearl, J. (2009). \textit{Causality: Models, Reasoning, and Inference} 
(2. Aufl.). Cambridge University Press.

\bibitem[van der Aalst(2016)]{vanderaalst2016process}
van der Aalst, W. M. P. (2016). \textit{Process Mining: Data Science in 
Action} (2. Aufl.). Springer.

\end{thebibliography}

\end{document}